\section{R-NEW003: a finite exact rational expression for the fourth series}
\label{sec:h4-finite-expression}

The expression below is a fully specified finite rational expression:
its numerator coefficients are finite integer character sums. It is not
an expanded minimal numerator, and no computation of all its coefficients
is asserted. It gives an unconditional finite reconstruction bound and
distinguishes that bound from experimental rational interpolation.

Let $\Phi_m$ denote the cyclotomic polynomial and define
\[
 D_4(t)=(1-t)^{19}\Phi_8(t)
       \prod_{\substack{2\le m\le24\\m\ne8,16}}\Phi_m(t)^{19}
       =\sum_{j=0}^{3196}d_jt^j.                    \tag{R-3a}
\]
Thus $d_0=1$ and $\deg D_4=3196$.
For $s\ge0$, and nonnegative integers $m_1,\ldots,m_{4s}$
satisfying $\sum_r rm_r=4s$, put
\[
\begin{split}
 Z(\mathbf m)&=\prod_{r=1}^{4s}r^{m_r}m_r!,\\
 C_s(\mathbf m)&=
 [x_1^{s+3}x_2^{s+2}x_3^{s+1}x_4^s]
 \prod_{1\le i<j\le4}(x_i-x_j)
 \prod_{r=1}^{4s}(x_1^r+x_2^r+x_3^r+x_4^r)^{m_r}.
\end{split}                                                   \tag{R-3b}
\]
Empty products give $Z(\varnothing)=C_0(\varnothing)=1$.
For $0\le j<3196$, define the explicit finite numbers
\[
 P_j=\sum_{\ell=0}^j d_\ell
       \sum_{\substack{m_1,\ldots,m_{4(j-\ell)}\ge0\\
                       \sum_r rm_r=4(j-\ell)}}
       \frac{C_{j-\ell}(\mathbf m)^3}{Z(\mathbf m)}.
                                                               \tag{R-3c}
\]

\begin{theorem}[R-NEW003]
The complete series in degree divided by four is
\[
 \boxed{\displaystyle H^{\rm ten}_4(t)
          =\frac{\sum_{j=0}^{3195}P_jt^j}{D_4(t)}.}             \tag{R-3d}
\]
Every $P_j$ is an integer. In particular, exact values of $k_4(s)$
for $0\le s\le3195$ determine the whole Hilbert series, with an
unconditional denominator and strict numerator-degree certificate.
The denominator displayed here is not asserted to be minimal.
\end{theorem}

\begin{proof}
The quotient has dimension $64-3(4^2-1)=19$: K-002--K-008
give an actual closed orbit with trivial effective stabilizer, so
the stable locus is nonempty and open. Hilbert--Serre rationality
places all poles at roots of unity. If $0<r<1$ and $|\xi|=1$,
nonnegative Hilbert coefficients give
\[
                  |H_4(r\xi)|\le H_4(r)=O((1-r)^{-19}).
\]
Hence every pole has order at most19. K-009 excludes roots of
order greater than $B_4=3125/128$, hence greater than24.
K-008 excludes primitive sixteenth roots and proves that the
primitive eighth-root poles are simple. Thus the actual reduced
denominator divides (R-3a), and $D_4H_4$ is a polynomial.
The degree is
\[
 19\sum_{m=1}^{24}\varphi(m)-19\varphi(16)-18\varphi(8)
       =19\cdot180-19\cdot8-18\cdot4=3196.
\]

The triple rectangular ray lies in the Kronecker cone, as already
witnessed by a critical tensor. Baldoni and Vergne (2015),
Proposition31, printed p.~23, states that its stretching function
is a quasipolynomial for every nonnegative integer, including zero;
their Theorem29 and the paragraph following it give the precise
Meinrenken--Sjamaar input. This is a cited theorem, not a novelty
claim. For completeness, the required properness consequence is
elementary. If $f(s)$ has period $p$ and constituent degree at most
$d$, each progression contributes
$t^aQ_a(t^p)/(1-t^p)^{d+1}$ with $\deg Q_a\le d$ and $0\le a<p$.
The summed numerator has degree at most $p(d+1)-1$, less than
the denominator degree. Hence $H_4$ is strictly proper, and
\[
                         \deg(D_4H_4)<3196.
\]
This implication is also recorded in Baldoni--Vergne, the
one-variable generating-series discussion on printed p.~20.

The ordinary Frobenius character formula identifies
$C_s(\mathbf m)$ with the symmetric-group character
$\chi^{(s,s,s,s)}$ on cycle type $(1^{m_1}2^{m_2}\cdots)$.
Indeed multiplication of the power-sum expansion by the four-variable
Vandermonde and extraction of the indicated strictly dominant
monomial selects precisely that Schur label. Character orthogonality
then gives
\[
                k_4(s)=\sum_{\sum rm_r=4s}
                       \frac{C_s(\mathbf m)^3}{Z(\mathbf m)}.
\]
These are the standard Frobenius and character-inner-product
identities (Macdonald, 1995, Chapter~I, equation~(7.3), pp.~113--114,
and equations~(7.5)--(7.6), p.~114), not new
coefficient formulas. Multiplication by $D_4$, followed by the
proved strict degree bound, gives (R-3c)--(R-3d). Integrality
follows because both $D_4$ and the Hilbert series have integral
coefficients.
\end{proof}

The file \texttt{certificates/H4\_denominator.json} records every
coefficient of the explicit denominator. The retained finite checks
through $s=12$ remain checks of the corresponding prefix, not a claim
that all3196 numerator coefficients were evaluated. This expression
does not yet compute $p(k_4)$ or $s(k_4)$, nor does it solve their
growth in variable dimension. Its new content is the explicit
finite bound obtained from K-008--K-009; the character identities
and quasipolynomiality theorem are classical inputs.
